\documentclass[preprint,sigplan]{acmart}

\usepackage{amsmath}
\usepackage{fancyvrb}
\usepackage{xypic}

\CustomVerbatimEnvironment{SVerbatim}{Verbatim}{fontsize=\small,
xleftmargin=0.5cm,
xrightmargin=0.2cm,commandchars=\\\{\},numbersep=0.9em}


%\newcommand{\cL}{{\cal L}}

\settopmatter{printacmref=false} % Removes citation information below abstract
\renewcommand\footnotetextcopyrightpermission[1]{} % removes footnote with conference information in first column

\newcommand{\dnote}[1]{\textcolor{purple}{Dom: #1}}
\newcommand{\tnote}[1]{\textcolor{blue}{Timo: #1}}

%% Bibliography style
\bibliographystyle{ACM-Reference-Format}
%% Citation style
%% Note: author/year citations are required for papers published as an
%% issue of PACMPL.
\citestyle{acmauthoryear}   %% For author/year citations

% Combining algegraic effects with monadfail
% Combining algegraic effects with pattern matching failure
% pattern matching failure due to representational issues without typing
%
\title{Algebraic Effects with Non-Recoverable Failure}

\author{Timotej Tomandl}

\affiliation{
%  \department{School of Computing}      %% \department is recommended
  \institution{University of Kent} %% \institution is required
%   \country{UK}                    %% \country is recommended
}

\author{Dominic Orchard}


\affiliation{
%  \department{School of Computing}      %% \department is recommended
  \institution{University of Kent} %% \institution is required
%   \country{UK}                    %% \country is recommended
}

\begin{document}

\pagestyle{plain} % removes running headers

\begin{abstract}
Algebraic effects are one of the possible approaches to tackle modularity issues arising in applications where normally we would use monad transformers.
One such area is programming language semantics, we show that the standard incarnation of effects does not satisfy some requirements needed to build compact,
non-repetitive semantics, when the language does not have a type system such as in cases of intermediate languages in compilers.
We show how to use a single MaybeT transformer to achieve this, we also explore the graded structure of effect monads.

\end{abstract}

\maketitle

% clearly seperate motivation
% failure and effects not satisfying MonadFail laws
% our construction, refer to Kiselyov
%    talk about MaybeT
%    gives the expected signature for handling
%    note this isnt exactly initiality
%
% illustrate the differences compared to Kiselyov's
% grading proofs
% short future work
% cleanup cakeml
% what matters first 2 pages can be longer
\section{Motivation}

Consider the following situation.  A compile writer is producing some
semantics-preserving translations between intermediate languages. At
some point in this transformation pipeline, information is lost, e.g.,
typing information is erased when converting to an untyped core
lanugage or byte code.  This necessitates that the compiled code has
some error handling cases, for example:
%
\begin{verbatim}
pseudo code
\end{verbatim}
%
This failure is intentionally non-recoverable. In Haskell, this
idea is captured by the \texttt{MonadFail} type class which is
used to desugar the \emph{do}-notation in the case of refutable
patterns. For example

\dnote{Show translation}

In this work, we seek to capture the same idea but in the
setting of \emph{algebraic effects}.

One aims of the algebraic-effects approach
is to allow the management of side effects to be
\emph{extensible}, considering first the necessary
effectful operations and their laws, rather than focussing on the underlying
semantic structure as is common in the \emph{monad}-first approach.
\emph{Monad transformers} can provide some modularity to the monad-first
approach, but this needs a new type class instance for every
combination of effects and this makes the implementation less
flexible.
\dnote{This is a very Haskell answer... is there something better here.}

% make clear our formalization verifies the laws but does not provide the syntactic do-solution


% flow of the ideas key properties basic structure

% work in progress the correct variant initiality
% normally the correct definition of handlers is defined in terms of initiality
% we would like to define not only handling, but handling in terms of algebra preservation
% we want maps to preserve equlities of presentations

% merge
% refer to constructions by reference
% explain differences between Kiselyov's, McBride's constructions
% ITrees, algebraic effects, McBride's Turing-completeness totally free
% Pirog, Gibbons Coinductive Resumption
% Plotkin Pretnar handling algebraic effects
% handlers of algebraic effects
% Grading should be a monoid after forgetting ordering
% Forgetful functor Monad -> Endo
% Forgetful functor Monad -> (ob(C),C)
% Failure is expressed by transformer construction which allows us to add an operation fail : M A with extra cancellative law
% failure is not initial as a forgetful functor (C,ob(C)) but with respect to MaybeT
% Free monad cannot satisfy the equation X >>= fail = X by itself because it would violate freeness as we have imposed an extra equation
% why not a failure with extra law, syntactic sugar
% can we get the syntactic sugar from quotients without imposing the MaybeT construction, no because handling as initiality is not well-behaved in this case, we want to really propagate failure
% further progress: algebraic freeness with HITs
% Continuations are not algebraic, delimited continuations should be thus providing a good way to work with exceptions
% Grading lifts through MaybeT, can we characterize other syntactic extensions of do-notation via transformers and laws, is there some generic quotient law for these cases
% characterization should work for all monad transformers which purely change return type
% what do we mean by grading lifting...
%

% running an effect can be done easily with an empty set because then the impure constructor can not hold any membership proof and thus is empty
%
%\section{Background}

\paragraph{Our results}

We have produced an intermediate construction between the variations of Kiselyov in Agda.
This construction is inbetween the construction the monads Eff and Freer as presented in their work. This is due to our focus on reasoning about monadic laws and
ease of construction within a dependently-typed framework instead of performance.
We further describe how to extend these with failing pattern matching, for simpler description, we present our code in stages and explain how to generically add
TODO add short explanation of similarities and differences
\dnote{Might be better to do the similarities and differences inline
  as we go.}

\section{Algebraic effects via graded monads}

Algebraic effects in theory consist of a signature of operations and
their equations.  An effect is represented by a type family handling
particular effect together with a name identifying the given effect.
% consists of a collection of operations and their signatures Ai -> Bi, these operations might
\dnote{Too Haskell.. can we say in a more general way.}

% signature consists of X, List (\a b : X -> a = b)
% and equations
% we do not impose any equations in our treatment of operations
% we verify monad and monadfail laws

\dnote{Oleg's definition}

The original definition of a Freer monad as defined by Kiselyov et al. is as follows
\begin{haskell}
data FFree f a where
  Pure :: a → FFree f  a
  Impure :: f x → (x → FFree f a) → FFree f a
\end{haskell}

\dnote{Then ours}
\begin{haskell}
data Eff (E : List (Set -> Set)) (A : Set) : Set₁ where
  Pure : A -> Eff E A
  Impure : {F : Set -> Set} {X : Set} -> .{proof_FinE : boolToSet (inSetList F E)} -> F X -> (X -> Eff E A) -> Eff E A
\end{haskell}


% running an effect
When the effect set is empty we should be able to extract the final
value by realizing that there can be no proof of membership within a
set.

%Monad transformers start with an encoding and only afterwards attempt to provide the necessary operations. Effects do the reverse we firt think in terms of operation and equations, which the
%y should satisfy and only provide implementation afterwards. An advantage of it is that the interpretation does not change structure of already written code, while this tends to happen when adding new operations using transformers. \cite{exteffold}


%\paragraph{Handlers}

%For example, this allows implementing stateful computation, which replies to a get ef%fect by returning the last put effects input.
%Algebraic effect handlers conform to a standard interface, which we call the `effect API', comprising the use of an effectful operation (called a `send', akin to throw for exceptions) and a way to interpret, or `handle', that operation (akin to catch for exceptions). We define later in more precise terms the form of the effect API, which are given in a generic way for any operation.
%Thus an effect API provides a way to `send' effects and `handle' them, where `send' corresponds to throwing and `handle' to catching an exception. It also provides a way to extract the final value once we have handled all effects.

\paragraph{Algebraic effects as monads and graded monads}

\dnote{define graded monads briefly, probably abstractly in mathematics}

Algebraic effects can be presented on top of a monadic interface, we provide this interface by implementing it using a free monad.
The importance of this interface is that we can provide an interface,
which chains different operations together and allows us to give a
correct interface.
\dnote{this sentence uses the word interface three times!}

For graded monads this allows us to do an interface, which combines
different effect sets.

\subsection{Effect monad construction}

We review the construction of Effect monad, to illustrate what issues pop up, while trying to make this into a admissible construction in Agda.
In haskell, these constructions can be done easily using the general datatypes, but they included datatype level fixed points in full.
In agda we can take 2 possible views of these constructions. We could go through the containers and construct a fixed point of containers as in \cite{?}.
Or we could fuse the left kan extension construction into the effect monad and observe that the result is strictly positive.
This means we do not have to provide a container construction for the underlying work as in \cite{?}. This works by realizing that the instantiation of F in the recursion to Lan of an arbitrary datatype is strictly positive without a recourse to containers.
We pick to do a halfway point between the constructions in Kyseliov, we do not use the optimized version, but we use the idea of effects should belong to a set.

\subsection{Graded monad proofs}
\dnote{is the intention here to explain about how we did the mechanisation?}
We grade monads to allow to combining the effect sets and thus being able to refer to only subcomputations in the programs and proofs.
We do this by describing a reindexing operation, which combines sets and allows us to refer to distinct effects with same names.
We need a way to reindex the effect set to embed the left effect set
into a list.
\dnote{the following sounds like it should go above in the definition
  of graded monads / the instance we are defining here.}
We note that the list of effects is itself a monoid and we grade by it. Graded unit, is the unit of our monad but for empty lists.
To produce graded bind, we append the 2 lists of effects and combine them. We


\paragraph{Handling effects}

%Handlers of effects can be thought of as a generalization of exception handlers, whic%h provide richer context
%and are constrained to respect a certain set of equations.
Algebraic effects
provide ways to write programs causing effects which are treated
abstractly initially but which are given concrete implementations
afterwards using \emph{handlers} (akin to raising an exception and
later handling it).
This generalizes to general effects: effect handlers provide a way to
not only handle general exceptions, but to also remember some state
between handling different operations a handler is responsible for.
Handlers are then also constrained to respect a set of domain-specific
equations.
\dnote{How much do we need to get into handlers here?}


\section{Leveraging MaybeT and MonadFail}

In Haskell, if a \emph{do}-expression performs a pattern
matching `bind' that may fail, then the monad being used
must also have an instance of the \texttt{MonadFail} class:
%
\begin{SVerbatim}
class MonadFail m where
  fail :: String -> m a
\end{SVerbatim}
%
where, for all monadic computations \texttt{x :: m a}:
%
{\small{
\begin{equation*}
\hspace{-10em}\texttt{fail $>$$>$$=$ x} \;\; \equiv \;\; \texttt{fail}
\end{equation*}
}}
%
i.e., failure is global and non-recoverable.

We given this a more abstract mathematical definition,
ignoring the \texttt{String} aspect in favour of a map
from a terminal object $1$.
%
\begin{definition}[Failing monad]
  For a category $\mathcal{C}$ with terminal object $1$
  then a \emph{failing monad} for endofunctor $T : \mathcal{C}
  \rightarrow \mathcal{C}$ is a monad with a natural transformation
  $\mathit{mfail}_A : 1 \rightarrow T\ A$ such that for all $g : 1
\rightarrow T\ B$ then:
$$
\mu \circ T (g \circ !_A) \circ \mathit{mfail}_A = \mathit{mfail}_B
$$
\end{definition}

If we attempt to give a \texttt{MonadFail} instance for the free
monad, we get stuck immediately: the free monad, is the monad for
which exactly the monad equations hold and nothing more, i.e.,
not the extra \texttt{MonadFail} axiom above. Therefore, we seek
to have both the benefits of the free monad structure for
capturing computation trees (for the purposes of effect handling)
but also the property of non-recoverable failure captured
by \texttt{MonadFail}. The simplest construction support the
above law generally is provided by the ``\texttt{MaybeT}''
\emph{monad transformer}~\cite{transformerpub}, which
in Haskell is typically defined as follows:
%
\begin{SVerbatim}
newtype MaybeT m a = MaybeT { runMaybeT :: m (Maybe a) }
\end{SVerbatim}
%
That is, a \texttt{MaybeT} is a monad homomorphism, mapping
$M -$ to $M (1 + -)$.
%
%MaybeT is the monad transformer, i.e. a transformation ${T A : Set} -> Monad T -> Maybe T A -> Monad S$.
%$MaybeT m x = m (Maybe x)$
%
We can think of \texttt{MaybeT} as an extension by an arbitrary,
single option failure case. Applying this to the free monad
structure yields a lawful instance of \texttt{MonadFail}:
%
\begin{SVerbatim}
  instance MonadFail (MaybeT m a) where
    fail _ = MaybeT (return Nothing)
\end{SVerbatim}
%
\dnote{verify equation?}
%
%especially in the strucutre of
%a Free monad as Monad is Pure and fixed point of a functor with this identity case.

However, we still require \emph{freeness} of this structure in order
to yield a suitable semantic core for handling algebraic effects
(augmented with non-recoverable failure). We therefore ask,
\emph{is the above the ``free MonadFail'' structure?}.

\dnote{merge in the next part}

\subsection{Handling EffectFail}
We have to modify handling, we state a different form of initiality, as much as effects are inspired by initiality, we provide a type-theoretic
motivation
for our definition, to provide handling of effectfail, we interpret the results and embed them into Maybe and propagate the underlying failure.
This should be correct not by some appeal to category theory as initiality is normally stated, but by appealing to the API we would expect from our system.
We would excpect this API to propagate pattern matching failure. To run EffectFail we either want to extract the underlying result or propagate the failure of underlying Pure result.

\paragraph{Algebraic freeness}

A monad $T$ is \emph{algebraically free on an endofunctor $F$} is
the category of $T$-algebras (monad algebras) is isomorphic to the category of
$F$-algebras. An algebraically-free monad is itself a free monad.
Therefore, in order to consider the notion of a ``free failing
monad'', we characterise the category of \texttt{MonadFail} algebras.

\begin{definition}
  For a monad $T : \mathcal{C} \rightarrow \mathcal{C}$, the category of monad algebras is the category
  of $T$-algebras which are compatible with the monad operations. That
  is $T\textsf{Alg}_{\mathsf{mnd}}$ is a category with objects as
  pairs $(A, h : A \rightarrow T\ A)$  where $A$ and $h$ are objects
  and morphisms of $\mathcal{C}$ respectively satisfying compatibility
  with the monad operations:
  %
  \dnote{Space for diagrams? or just summarise as equation}
\end{definition}

\begin{definition}
  For a category $\mathcal{C}$ with terminal object $1$ and
  a failing monad $T : \mathcal{C} \rightarrow \mathcal{C}$, the
  category of its failing monad algebras
  $T\textsf{Alg}_{\mathsf{mndFail}}$ is a category of monad algebras, with objects
  pairs $(A, h : T\ A \rightarrow A)$  where $A \in |\mathcal{C}|$ and
  $h \in \mathcal{C}(T\ A, A)$ compatible with the monad operations,
  and the following additional compatibility with $\mathit{mfail}$ if
  there exists a morphism $h' : 1 \rightarrow A$ such that:
  %
  \begin{equation*}
    \xymatrix@C=4em{
      1 \ar[dr]_{h'} \ar[r]^-{\mathit{mfail}_A} &  T\ A \ar[d]^{h} \\
      & A
      }
  \end{equation*}
  %
\end{definition}
%
\dnote{Result that the free monad algebra with $(T\ A, \mu)$ is also
  the free failing monad algebra where the above equation comes from
  the \texttt{MonadFail} axiom.}


\section{Future directions}
As much we have constructed the monad in agda and verified the various laws. For practical purposes there needs to be more than just a monadic structure.
We want to reason about the local code using the eponymous algebraic laws. We would like to attach equations the handlers should satisfy to the monad.
There are multiple interesting directions to take this in.
We would also like to use Quotient inductive types to provide a view of operations for reasoning and correct-by-construction definition of handlers,
which abstracts away from the concrete choices when reasoning, yet forces us to prove satisfiability of the algebriac equations by some handler.
To extend the proofs from our current work to the quotient case we need to prove a type-theoretic formulation of algebraic freedom \cite{ncatlab}.
To verify that all the constructions preserve all equations. We note that by theorem 3.1 and 3.2 \cite{ncatlab} this result should follow in Homotopy type theory.
The reason for this being that we need the ambient category to be locally small and complete with homotopy type theory we should be able to get all colimits.

These lemmas apply only to the free monad with a forgetful functor $Mnd(C) \rightarrow Endo(C)$, we suggest that this definition should extend to $Mnd(C) \rightarrow (C,ob(C))$.
We conjecture that the F-algebra of the left kan extension endofunctor for a quotient inductive type family is the underlying algebra as in the initial algebra semantics of the given type family.
This should allow us to extend the isomorphism in terms of algebraic freeness to the freer construction.
We would then expect the freer monad T over X to have T-algebras isomorphic to the algebras of X.
We would expect the formal statement of this conjecture might be provable in some appropriate container \cite{indexed containers} or generic levitated presentation \cite{levitation},
which would allow us to connect the F-algebra of a functor with the algebra of the underlying datatype.

Another direction of future research is to construct an isomorphism between continuations and the fast type-aligned queues of Kiselyov, quotiented by their composition,
thus providing a more efficient implementation of effects without the overheads to reasoning this entails as we focus on semantic constructions instead of executable code, this is not our primary concern.
We would expect this to then transport by univalence, but as discussed
in {unitransport} we would need also some form of representation
independence coming from parametricity for this to be actually
efficient.

\appendix
\section{Appendix Title}

This is the text of the appendix, if you need one.

% The 'abbrvnat' bibliography style is recommended.

\bibliography{references}

\end{document}
